\chapter{От требований физики к планковскому миру}
\label{ch:descent}

Закон, записанный мельче уже проверенного масштаба, не свободен.
После осреднения он обязан вернуть те уравнения и те запреты, которые прибор уже сверил.
Ниже мы спускаемся по этажам и на каждом записываем, во что это обязательство превращается.
Планковский мир в конце есть то место, где непрерывное поле в точке перестаёт с обязательством справляться.

\section{Требования физики}
\label{sec:requirements}

Обязательство одно: нижний этаж воспроизводит верхний.
Чтобы было с чем сверять, скажем, что на макроскопическом масштабе установлено прямо, без гипотезы об устройстве вещества.

Возмущение доходит за конечное время.
Если бы влияние было мгновенным, показание прибора в данном месте зависело бы от того,
что в тот же момент происходит сколь угодно далеко, и локальный опыт не сходился бы.
За время~$\Delta t$ сигнал проходит путь не длиннее~$c\Delta t$:
\begin{equation}
  |\Delta\mathbf{x}| \le c\,\Delta t.
  \label{eq:light-cone}
\end{equation}
Любой более глубокий закон обязан сохранить этот конус.

Из \eqref{eq:light-cone} следует локальность.
Далёкое не успевает прийти, поэтому новое состояние в данной точке собирается из поля в её окрестности.
Для электромагнитного поля это уравнения Максвелла:
производная по времени стоит в точке, ротор и дивергенция берут соседние значения,
\begin{align}
  \nabla\times\mathbf{E} &= -\partial_t\mathbf{B}, &
  \nabla\cdot\mathbf{B} &= 0, \label{eq:maxwell-hom}\\
  \nabla\times\mathbf{B}
    &= \mu_0\mathbf{j} + \mu_0\varepsilon_0\,\partial_t\mathbf{E}, &
  \nabla\cdot\mathbf{E} &= \rho/\varepsilon_0. \label{eq:maxwell-inh}
\end{align}
Характеристики этой системы лежат на конусе \eqref{eq:light-cone}.
Та же схема у гидродинамики.
Для идеальной жидкости скорость в точке меняется от давления и скорости рядом:
\begin{equation}
  \partial_t\mathbf{v} + (\mathbf{v}\cdot\nabla)\mathbf{v}
  = -\frac{1}{\rho}\nabla p.
  \label{eq:euler}
\end{equation}
В механике сплошной среды справа стоит дивергенция напряжений,
$\partial_k\sigma_{ik}$, снова из окрестности той же точки.

Механическое состояние исчерпывается координатами и скоростями.
Положение точки есть $\mathbf{r}(t)$, скорость --- его производная,
\begin{equation}
  \mathbf{v} = \dot{\mathbf{r}}.
  \label{eq:velocity}
\end{equation}
Уравнения Ньютона связывают ускорение с этим состоянием:
\begin{equation}
  m\ddot{\mathbf{r}} = \mathbf{F}(\mathbf{r},\dot{\mathbf{r}}).
  \label{eq:newton}
\end{equation}
Это уравнение второго порядка.
Задание $\mathbf{r}$ и~$\dot{\mathbf{r}}$ в один момент восстанавливает их и в предыдущий, и в следующий.
Старшей производной среди начальных данных в \eqref{eq:newton} нет:
закон, которому она понадобилась бы, описывал бы другой мир.

Тепловое равновесие добавляет два запрета, тоже измеренных.
В адиабатически замкнутой системе энтропия не убывает,
\begin{equation}
  \mathrm{d}S \ge 0.
  \label{eq:second-law}
\end{equation}
При $T\to 0$ энтропия равновесной системы перестаёт зависеть от параметров равновесия~$\lambda$
и в шкале Планка обращается в нуль:
\begin{equation}
  \Bigl(\frac{\partial S}{\partial\lambda}\Bigr)_{T}
    \xrightarrow[T\to 0]{} 0,
  \qquad
  S \xrightarrow[T\to 0]{} 0.
  \label{eq:third-law}
\end{equation}
Система садится в одно основное состояние.
Закон, у которого при нулевой температуре оставалось бы макроскопически много
равноправных состояний, \eqref{eq:third-law} противоречит.

Вакуум, насколько его видит свет, не выделяет направления:
скорость одна и та же вдоль любой оси.
Вместе с \eqref{eq:light-cone} это интервал
\begin{equation}
  \mathrm{d}s^2 = c^2\,\mathrm{d}t^2 - \mathrm{d}\mathbf{x}^2,
  \label{eq:interval}
\end{equation}
один и тот же во всех инерциальных системах.
Заряд в замкнутой системе сохраняется.
Это уравнение непрерывности
\begin{equation}
  \partial_t\rho + \nabla\cdot\mathbf{j} = 0,
  \label{eq:continuity}
\end{equation}
а на узле проводников --- закон Кирхгофа, $\sum_k I_k = 0$.
Не свойство, приписанное заранее микроскопическому полю.

На этом макроскопическая сверка кончается.
Спин, комплексная амплитуда, целочисленность устойчивого заряда
и вопрос о том, откуда берутся безразмерные числа, здесь ещё не видны.
Они появятся на том этаже, где без них уже проверенный опыт перестаёт получаться,
и каждое окажется тем же конусом, тем же сохранением или той же обратимостью, дочитанными до этого этажа.

\section{Макромир}
\label{sec:macro-known}

Макромир есть эти же уравнения, записанные для непрерывных полей.

Из \eqref{eq:second-law} и \eqref{eq:third-law} строится термодинамика.
Температура есть производная энергии по энтропии,
\begin{equation}
  T = \Bigl(\frac{\partial E}{\partial S}\Bigr)_{V,N},
  \label{eq:temperature}
\end{equation}
и из неё же --- потенциалы.
Интервал \eqref{eq:interval} один и тот же во всех инерциальных системах,
поэтому энергия и импульс составляют четырёхвектор
\begin{equation}
  p^{\mu} = (E/c,\,\mathbf{p}).
  \label{eq:four-momentum}
\end{equation}

Из этих уравнений собрана система единиц.
Числа~$c$, $\hbar$, $k_B$, заряд~$e$, массы частиц и постоянная тонкой структуры~$\alpha_{\mathrm{fs}}$ в ней измерены.
Сами уравнения их не выводят.
На этом этаже такое ещё терпимо: число может быть константой прибора.
Один этаж ниже тот же закон, если он претендует быть фундаментальным, уже не имеет права оставлять число вписанным рукой.

\section{Микромир}
\label{sec:macro-open}
\label{sec:microworld}

Внутри вещества непрерывные уравнения, только что записанные, ещё держат механику и электродинамику прибора.
Они перестают держать сам опыт, как только отсчёт становится единичным.

Ансамбль подготовленных одинаково атомов даёт частоты исходов.
Один запуск прибора даёт один исход.
Уравнения поля содержат частоты и не содержат границу между системой, прибором и записью результата.
Сохранение того, что в ансамбле выглядит как полная вероятность, здесь ещё не объясняет одиночный отсчёт.

Числа~$c$, $\hbar$, $G$, $e$ и массы фермионов измерены с большой точностью.
Безразмерное число, которое из них собирается, --- постоянная тонкой структуры
\begin{equation}
  \alpha_{\mathrm{fs}} = \frac{e^2}{4\pi\varepsilon_0\hbar c}.
  \label{eq:alpha-fs}
\end{equation}
Она стоит в законе свободным множителем:
измерение есть, вывода из \eqref{eq:maxwell-hom}--\eqref{eq:newton} нет.
Пока закон претендует только описать прибор, такой множитель терпим.
Как только тот же закон считается фундаментальным, вписанное число означает, что измерение не выведено.

Чёрная дыра делает ограниченность плотности видимой числом.
Энтропия горизонта пропорциональна его площади~$A$:
\begin{equation}
  S = \frac{k_B A}{4\ell_P^2},
  \qquad
  \ell_P^2 = \frac{\hbar G}{c^3}.
  \label{eq:bh-entropy}
\end{equation}
Число различимых состояний в причинно связанной области поэтому конечно,
\begin{equation}
  \ln\mathcal{N} \le \frac{A}{4\ell_P^2}.
  \label{eq:state-count}
\end{equation}
У шара радиуса~$R$ правая часть пропорциональна~$R^2$.
Набор независимых значений поля в каждой точке объёма рос бы как объём, то есть как~$R^3$,
а непрерывное значение в каждой точке $x\in\mathbb{R}^3$ даёт несчётное множество конфигураций.
Энергия в той же области ограничена сверху по той же причине.
Горизонт Шварцшильда
\begin{equation}
  R_s = \frac{2GE}{c^4}
  \label{eq:schwarzschild}
\end{equation}
достигает радиуса области, когда энергия слишком велика:
\begin{equation}
  E \le \frac{R c^4}{2G}.
  \label{eq:energy-bound}
\end{equation}
Общая теория относительности и квантовая теория поля локально друг другу не противоречат
и счёта \eqref{eq:state-count} не содержат.

Там, где одновременно существенны $\hbar$, $G$ и~$c$, масштаб есть~$\ell_P$ из \eqref{eq:bh-entropy}.
Поле в каждой точке счёт \eqref{eq:state-count} нарушает.
Следующий этаж --- запись того же сохранения и той же обратимости,
при которой спектр атома и два знака спина уже получаются.

\section{Квантовая механика}
\label{sec:qm}

Классическая траектория не выдерживает атом:
устойчивые орбиты и линейчатый спектр ей противоречат.
Кулоновская задача для электрона,
\begin{equation}
  \Bigl(-\frac{\hbar^2}{2m_e}\nabla^2 - \frac{e^2}{4\pi\varepsilon_0 r}\Bigr)\psi = E\psi,
  \label{eq:hydrogen}
\end{equation}
даёт дискретные уровни
\begin{equation}
  E_n = -\frac{m_e c^2\alpha_{\mathrm{fs}}^2}{2n^2},
  \qquad n=1,2,\ldots
  \label{eq:rydberg}
\end{equation}
Опыт интерференции показывает ещё и то, что складываются не частоты исходов, а амплитуды:
\begin{equation}
  w \propto |\psi_1+\psi_2|^2,
  \label{eq:interference}
\end{equation}
иначе тёмных полос не было бы.
Состояние, из которого такие амплитуды читаются, есть вектор;
частота исхода на макроэтаже есть квадрат его модуля,
\begin{equation}
  w(a) = |\langle a|\psi\rangle|^2.
  \label{eq:born}
\end{equation}
Обратимость, уже видная в \eqref{eq:newton}, здесь означает,
что эволюция не меняет полную норму.
Уравнение Шрёдингера
\begin{equation}
  i\hbar\,\partial_t\psi = H\psi,
  \qquad H^\dagger = H,
  \label{eq:schrodinger}
\end{equation}
сохраняет
\begin{equation}
  \langle\psi|\psi\rangle = \mathrm{const}
  \label{eq:norm}
\end{equation}
и по состоянию в один момент восстанавливает состояние в предыдущий.

Опыт со спином электрона даёт два исхода,
\begin{equation}
  s_z = \pm\frac{\hbar}{2},
  \label{eq:spin-z}
\end{equation}
и ещё одно свойство поворота.
Поворот на угол~$\theta$ вокруг оси~$\mathbf{n}$ действует как
\begin{equation}
  U(\mathbf{n},\theta)
  = \exp\bigl(-i\theta\,\mathbf{n}\cdot\boldsymbol{\sigma}/2\bigr).
  \label{eq:spin-rotation}
\end{equation}
Оборот на $2\pi$ меняет знак, оборот на $4\pi$ возвращает состояние:
\begin{equation}
  U(2\pi)=-\mathbb{1},
  \qquad
  U(4\pi)=+\mathbb{1}.
  \label{eq:spin-4pi}
\end{equation}
Запись состояния --- столбец $\psi\in\mathbb{C}^2$.
Общий множитель-фаза $e^{i\alpha}$ не меняет частот \eqref{eq:born}
и соответствует сохранению заряда, уже записанному \eqref{eq:continuity}.
Локальный поворот двух компонент лежит в $SU(2)$.

Физическая величина, построенная из такого состояния,
не имеет права зависеть от того, где проведён разрез аргумента на $2\pi$.
Комбинация, которая разреза не видит, собирается из амплитуды и сопряжённой амплитуды:
\begin{equation}
  \langle\psi|O|\psi\rangle.
  \label{eq:observable}
\end{equation}

Квантовая механика закрывает эти факты для одной системы с заданным числом степеней свободы.
Одиночный отсчёт прибора по-прежнему стоит рядом с уравнением.
Когда степеней свободы много и частицы рождаются и исчезают,
та же норма и тот же световой конус записываются как теория поля.

\section{Теория поля}
\label{sec:qft-open}

На масштабах ядра и элементарных частиц тот же конус, та же локальность и тот же спин
читаются как лагранжиан с калибровочными симметриями $SU(3)\times SU(2)\times U(1)$,
фермионами поколений и механизмом Хиггса.
Квантовая электродинамика и хромодинамика дают проверенные поправки.

Заряд, который доходит до прибора, кратен элементарному:
\begin{equation}
  Q = n e, \qquad n\in\mathbb{Z}.
  \label{eq:charge-quantum}
\end{equation}
Свободный кварк не наблюдается, заряд лептона целый.
Закон, у которого устойчивый дефект нёс бы произвольную долю этого кванта,
таблицу частиц не воспроизводит.
Произведение зеркального отражения, зарядового сопряжения и обращения времени опыт сохраняет,
\begin{equation}
  (\mathrm{CPT})\,\mathcal{L}\,(\mathrm{CPT})^{-1} = \mathcal{L};
  \label{eq:cpt}
\end{equation}
по отдельности отражение и зарядовое сопряжение в слабом взаимодействии нарушены.
Общий поворот фазы вакуума $\psi\mapsto e^{i\alpha}\psi$ не меняет частот:
это то же сохранение заряда, \eqref{eq:continuity}.
Обе симметрии должны сидеть в самом законе.

Отсюда же видно, где непрерывная запись перестаёт справляться.

Возмущёние в квантовой теории поля содержит расходимости.
Группа ренормализации снимает их с наблюдаемых величин
и оставляет вопрос, что стоит за взаимодействием там,
где оно перестаёт быть точечным.
Это то же ограничение плотности, теперь внутри возмущения.

Стандартная модель содержит более девятнадцати настраиваемых величин:
массы, углы смешивания, константы связи.
Симметрии $SU(3)\times SU(2)\times U(1)$ их не фиксируют.
Как и этажом выше, числа стоят в законе вписанными.

Квантование метрики в той же схеме не замкнуто.
Планковская энергия
\begin{equation}
  E_P = \sqrt{\frac{\hbar c^5}{G}}
  \label{eq:planck-energy}
\end{equation}
--- граница записи тех же требований.
Теория поля --- последний непрерывный этаж.

\section{Планковский мир}
\label{sec:planck-floor}
\label{ch:carrier-object}
\label{sec:carrier-object}

Масштаб длины, на котором одновременно видны
квантовость, локальная кривизна и причинность, уже записан: $\ell_P$ в \eqref{eq:bh-entropy}.
Вместе с ним время~$t_P=\ell_P/c$ и масса~$m_P=\sqrt{\hbar c/G}$,
а энергия того же масштаба --- \eqref{eq:planck-energy}.
Физика по-прежнему не зависит от человеческих единиц:
важны безразмерные отношения и те же симметрии, что названы выше.

Требование ограниченной плотности уже названо.
Если в каждой точке $x\in\mathbb{R}^3$ сидит своё комплексное число,
в любом конечном причинном объёме конфигураций несчётно много,
а требование уже ограничивает это число.
Место, на котором это выполнено, --- счётное множество узлов.

\begin{definition}[Носитель как множество узлов]
\label{def:carrier-sites}
Носитель микроуровня~$M$ --- счётное множество узлов~$\Lambda$.
Время, согласное с одним шагом сигнала, дискретно: $t\in\mathbb{Z}$.
\end{definition}

Требование конечной скорости уже названо.
На счётном множестве узлов конечное расстояние за конечное время --- это один шаг за один такт.

\begin{definition}[Элементарные шаг и такт]
\label{def:step-tick}
$\hl$ --- элементарный шаг: длина, на которую за один такт уходит сигнал.
$\htick$ --- один такт эволюции.
Тактовая скорость сигнала на~$M$
\[
  c_0 = \frac{\hl}{\htick}.
\]
\end{definition}

\begin{definition}[Окрестность одного такта]
\label{def:neighborhood}
Окрестность узла $x\in\Lambda$ --- множество узлов, достижимых из~$x$
ровно за один такт:
\[
  N(x)\subset\Lambda.
\]
\end{definition}

Требования спина и сохранения нормы уже записаны в квантовой механике.
На узле они остаются теми же:
значение --- спинор из двух компонент, норма --- сумма квадратов модулей.

\begin{definition}[Поле и норма]
\label{def:field-a3}
На каждом узле $x\in\Lambda$ задан спинор
\[
  z(x,t)=(z_1,z_2)^\top\in\mathbb{C}^2.
\]
$\PlanckCell$ --- этот узел.
Модуль $|z_1|^2+|z_2|^2$ --- плотность поля на узле;
фазы $z_k/|z_k|$ --- степени свободы, которые перемешиваются при эволюции.
Конфигурация всех узлов в такт~$t$ есть
\[
  \Psi^t=\{z(x,t)\}_{x\in\Lambda}.
\]
Суммарная норма, аналог $\langle\psi|\psi\rangle$,
\[
  N(t):=\sum_{x\in\Lambda} z^\dagger(x,t)\,z(x,t)=\sum_x\bigl(|z_1|^2+|z_2|^2\bigr).
\]
\end{definition}

Дальше --- не новый «язык'', а те же требования, записанные без континуума:
дискретная фаза, конечный регистр на узле, обратимый локальный шаг.
Геометрия~$\Lambda$ (FCC) и конкретный закон~$g$ выбираются в следующих главах;
здесь фиксируется, \emph{что} осталось после QED и планковского счёта.

\section{Дискретизация фазы}
\label{sec:descent-phase-disc}
\label{sec:coords-internal}

На макроэтаже волновая функция допускает запись Маделунга,
\[
  \psi=\sqrt{\rho}\,e^{i\varphi/\hbar},
\]
с плотностью $\rho=|\psi|^2$ и фазой~$\varphi$.
На узле~$\PlanckCell$ тот же смысл несут $|z|^2$ и аргументы компонент $z_k$.
Континуумная фаза $\varphi\in\mathbb{R}$ на одной планковской ячейке не определена бесконечно точно:
требование обратимости и конечного числа различимых состояний
(\eqref{eq:state-count}) задаёт \emph{конечное кольцо фаз}.

Счёт \eqref{eq:state-count} на ячейке объёма~$v_{\hv}\sim\hl^3$ даёт информационный бюджет
\[
  B_{\hv}=\frac{2\pi}{\ln 2}\quad\text{(мост $T$)},\qquad
  N_{\mathrm{ring}}=2^{\lfloor B_{\hv}\rfloor}=512,\qquad
  B_{\hv}=\log_2 N_{\mathrm{ring}}+\varepsilon_B,\quad \varepsilon_B\approx 0{,}065.
\]
Минимальный фазовый шаг, уже видимый в квантовой механике как $\Delta\phi_{\min}=\tfrac12$,
на кольце есть целое число тиков:
\[
  \Delta\phi_{\mathrm{disc}}=41,\qquad
  N_\phi=\lceil 2\pi/\Delta\phi_{\min}\rceil=13.
\]
Один такт~$\htick$ --- полный оборот $\mathbb{Z}_{N_{\mathrm{ring}}}$;
генератор вращения $\omega=e^{i2\pi/N_{\mathrm{ring}}}$, состояние на $S^1$ --- $\omega^{\varphi_{\mathrm{disc}}}$.

\begin{remark}[Фаза, $\omega$ и kick $\Phi$]
\label{rem:phi-omega-Phi}
$\varphi_{\mathrm{disc}}\in\mathbb{Z}_{N_{\mathrm{ring}}}$ --- \emph{координата} фазы (аддитивный индекс).
$\omega$ --- генератор, не вторая координата.
$\Phi\in\mathbb{Z}_{N_{\mathrm{ring}}}$ --- фазовый \emph{kick} за такт ($R(\Phi)=\omega^{\Phi}$), не регистр $\varphi_{\mathrm{disc}}$.
В мосте~$T$: $\varphi=\varphi_{\mathrm{disc}}\cdot(2\pi/N_{\mathrm{ring}})$;
поправка Швингера $a_e\sim\alpha_{\mathrm{fs}}/(2\pi)$ читается там же, на~$M$ --- в тиках кольца.
\end{remark}

Спинор на узле кодируется в конечном кольце Гауссовых целых
\[
  Z=U+iV\in\mathbb{Z}_{N_{\mathrm{ring}}}[i],
\]
амплитуда --- на сетке $Q(2^{-6})$, $\mu\in\mathbb{Z}_{64}$.
Дискретная фаза раскладывается по Heisenberg-кванту:
\[
  \varphi_{\mathrm{disc}} \equiv k_\phi\cdot\Delta\phi_{\mathrm{disc}}+\varphi_{\mathrm{f}}
  \pmod{N_{\mathrm{ring}}},
\]
$k_\phi\in\mathbb{Z}_{13}$ --- сектор, $\varphi_{\mathrm{f}}\in\{0,\ldots,40\}$ --- тонкая фаза.

\begin{proposition}[Каноническая фаза и шов]
\label{prop:canonical-phase}
Каждому $\varphi_{\mathrm{disc}}\in\mathbb{Z}_{512}$ сопоставляется единственная каноническая пара
$(k_\phi,\varphi_{\mathrm{f}})$ с минимальным $k_\phi$.
Наивная карта $\mathbb{Z}_{13}\times\mathbb{Z}_{41}\to\mathbb{Z}_{512}$ имеет $21$ коллизию;
\[
  N_{\mathrm{ring}}=N_\phi\cdot\Delta\phi_{\mathrm{disc}}-(N_\phi+N_{\mathrm{hier}})=13\cdot 41-21.
\]
\end{proposition}

Фазовый остаток планковской дырки на~$M$: $\alpha_*-1=\Delta\phi_{\mathrm{disc}}/N_{\mathrm{ring}}=41/512$;
в мосте~$T$ то же число носит имя $1/(4\pi)$ и не совпадает с $\alpha_{\mathrm{fs}}$ из \eqref{eq:alpha-fs}.

\section{Обобщённые координаты}
\label{sec:coords-dh-dt}
\label{ch:native-coords}

\begin{definition}[Внешние координаты $(dh,dt)$]
\label{def:dh-dt}
\[
  dh:=\hl,\qquad dt:=\htick,\qquad c_0=dh/dt.
\]
В натуральных единицах симуляции $dh=dt=1$.
Одна ступень действия~$s_0$ на $(\hv,1\,dt)$ порождает лестницу
$p_0=s_0/dh$, $E_0=s_0/dt$, $F_0=s_0/(dh\cdot dt)$ без трёх независимых подгонок.
\end{definition}

\begin{definition}[Внутренние координаты]
\label{def:internal-coords}
Регистр ячейки: $(\mu,k_\phi,\varphi_{\mathrm{f}})$ или $(\mu,\varphi_{\mathrm{disc}})$;
спинор~$\mathbb{C}^2$ --- две копии; вакуум: $\varphi_{\mathrm{f}}=0$, $\varphi_{\mathrm{disc}}=k_\phi\cdot 41$.
\end{definition}

\begin{table}[ht]
\centering
\caption{Лестница $s_0$ и мост $M\to T$.}
\label{tab:dh-dt-bridge}
\begin{tabular}{@{}lll@{}}
\toprule
На $M$ & $(dh,dt)$ & Мост $T$ \\
\midrule
$s_0$ & на $(\hv,1\,dt)$ & $\hbar/2$ \\
$p_0$ & $s_0/dh$ & $\hbar/(2\hl)$ \\
$E_0$ & $s_0/dt$ & $\hbar/(2\htick)$ \\
$c_0$ & $dh/dt$ & $\hl/\htick$ \\
\bottomrule
\end{tabular}
\end{table}

\section{Локальный обратимый автомат}
\label{sec:descent-ca}

Когда на каждом узле конечный алфавит, время дискретно, а обновление локально,
микроуровень~$M$ есть \emph{клеточный автомат}:
\[
  s(x,t+1)=f\bigl(\{s(y,t)\}_{y\in N(x)}\bigr),
\]
но с физическими ограничениями из спуска: сохранение нормы, обратимость, спин~$1/2$, заряд~$\in\mathbb{Z}$.

\begin{definition}[Закон эволюции]
\label{def:g-ca}
Глобальная конфигурация $\Psi^t=\{z(x,t)\}_{x\in\Lambda}$ обновляется локальным обратимым правилом
\[
  \Psi^{t+1}=g(\Psi^t),\qquad g^{-1}\circ g=\mathrm{id},
\]
с конечным регистром $Z\in\mathbb{Z}_{N_{\mathrm{ring}}}[i]$ на узле.
Исторические примеры (Конвей, Rule~110) иллюстрируют класс CA, но не задают~$g$;
см. приложение~\ref{app:background}.
\end{definition}

\subsection{Словарь: термины клеточного автомата и слой~\texorpdfstring{$M$}{M}}
\label{sec:ca-terminology}

В литературе по клеточным автоматам и «цифровой физике'' те же объекты часто называют иначе,
чем в QED и в гл.~\ref{ch:carrier} о двух скоростях.
Ниже --- соответствие без отождествления микро- и макроскорости.

\begin{table}[ht]
\centering
\caption{КА-терминология и наша запись на~$M$.}
\label{tab:ca-glossary}
\small
\begin{tabular}{@{}p{3.4cm}p{4.8cm}p{4.6cm}@{}}
\toprule
Термин КА & На~$M$ & Замечание \\
\midrule
ячейка, site & узел $\hv\in\Lambda$ & регистр $Z\in\mathbb{Z}_{N_{\mathrm{ring}}}[i]$ \\
состояние $s(x)$ & спинор $z(x)\in\mathbb{C}^2$ & в коде --- фиксированная точка на $Z$ \\
окрестность & $N(x)$ (\cref{def:neighborhood}) & световая за один такт, не Мур \\
правило, $f$ & закон $g$ (\cref{def:g-ca}) & обратим; $g^{-1}\circ g=\mathrm{id}$ \\
поколение, generation & такт $dt=\htick$ & $t\in\mathbb{Z}$ \\
шаг решётки & $dh=\hl$ & один hop по $\Lambda$ \\
\emph{скорость света} CA & $c_0=dh/dt=\hl/\htick$ & \textbf{тактовая} скорость сигнала \\
макроскорость & $c=\kgeom c_0$ & мост~$T$, CODATA (\cref{sec:two-c}) \\
световой конус & $|x-y|\le c_0\cdot\Delta t\cdot dh$ & каузальность одного такта \\
алфавит & $\mathbb{Z}_{N_{\mathrm{ring}}}[i]\times Q(2^{-6})$ & конечный, не $\{0,1\}$ \\
\bottomrule
\end{tabular}
\end{table}

В учебниках по CA «скорость света'' почти всегда означает \emph{не} экспериментальную $c$,
а нормировку «одна ячейка за один такт'': $c_0=1$ в natural units.
У нас то же: за один~$dt$ сигнал не дальше одного ребра длины~$dh$ (\cref{def:step-tick}).
Макроскопическая скорость света в вакууме~$c$ появляется только после осреднения на~$T$
и связана с~$c_0$ коэффициентом~$\kgeom$ (гл.~\ref{ch:carrier}, \cref{sec:two-c}).
Путать $c_0$ и $c$ --- типичная ошибка при чтении CA-литературы рядом с физикой поля.

Закон~$g$ не подбирается под таблицу Rule~110: он \emph{замыкается} из абсолютных условий
и геометрии носителя (гл.~\ref{ch:axioms}, \ref{ch:carrier}, \ref{ch:evolution}).
Макромир~$T$, с которого начался спуск, --- осреднение $\mathcal{B}$ по~$M$;
гладкие уравнения QED и Маделунга --- его имена, не ядро автомата.

\section{Сводка моста}
\label{sec:coords-bridge}

\begin{table}[ht]
\centering
\caption{Где живут $\pi$, $\ln 2$ и CODATA.}
\label{tab:m-t-bridge}
\begin{tabular}{@{}p{0.42\linewidth}p{0.48\linewidth}@{}}
\toprule
Ядро $M$ & Только мост $T$ \\
\midrule
$\varphi_{\mathrm{disc}}$, $k_\phi$, $\mu$; $dh$, $dt$ & $\hbar$, $t_P$, $\ell_P$, $c$ \\
$\Delta\phi_{\mathrm{disc}}/N_{\mathrm{ring}}$ & $\alpha_*-1=1/(4\pi)$ \\
$\log_2 N_{\mathrm{ring}}+\varepsilon_B$ & $B_{\hv}=2\pi/\ln 2$ \\
$N_\phi=13$ & $\lceil 4\pi\rceil$ \\
\bottomrule
\end{tabular}
\end{table}

Спуск с макроскопической физики на~$T$ к обратимому автомату на~$M$ завершён.
Далее: выбор геометрии~$\Lambda$ (FCC), теорема дискретности и явный закон~$g$.
